\section{Problem Set 3: Matrix Inversion}

\begin{enumerate}
\item \textbf{Problem 1}\par
Decide whether each matrix is invertible:
\[
A=\begin{pmatrix}2&1\\3&4\end{pmatrix},
\qquad
B=\begin{pmatrix}1&2\\2&4\end{pmatrix}.
\]

\item \textbf{Problem 2}\par
Find the inverse of
\[
A=\begin{pmatrix}3&1\\2&1\end{pmatrix}
\]
using the \(2\times2\) inverse formula.

\item \textbf{Problem 3}\par
Verify your answer to the previous problem by computing both \(AA^{-1}\) and \(A^{-1}A\).

\item \textbf{Problem 4}\par
Find the inverse of
\[
D=\operatorname{diag}(2,-3,5).
\]

\item \textbf{Problem 5}\par
Find all values of \(k\) for which
\[
A(k)=\begin{pmatrix}
k&1\\
2&k
\end{pmatrix}
\]
has an inverse.

\item \textbf{Problem 6}\par
Use Gauss--Jordan elimination to find the inverse of
\[
A=\begin{pmatrix}
1&2\\
3&5
\end{pmatrix}.
\]

\item \textbf{Problem 7}\par
Use Gauss--Jordan elimination to find the inverse of
\[
A=\begin{pmatrix}
1&1&0\\
0&1&1\\
1&0&1
\end{pmatrix}.
\]

\item \textbf{Problem 8}\par
Find the adjugate of
\[
A=\begin{pmatrix}
4&1\\
2&3
\end{pmatrix}
\]
and use it to compute \(A^{-1}\).

\item \textbf{Problem 9}\par
Let
\[
A=\begin{pmatrix}2&1\\1&1\end{pmatrix},
\qquad
b=\begin{pmatrix}5\\3\end{pmatrix}.
\]
Solve \(Ax=b\) using \(A^{-1}\).

\item \textbf{Problem 10}\par
Solve
\[
XA=B
\]
for \(X\), where
\[
A=\begin{pmatrix}1&1\\0&2\end{pmatrix},
\qquad
B=\begin{pmatrix}3&1\\4&2\end{pmatrix}.
\]
Pay attention to the side on which \(A^{-1}\) must be multiplied.

\item \textbf{Problem 11}\par
Solve
\[
2X-A=B
\]
for \(X\), where
\[
A=\begin{pmatrix}1&0\\2&-1\end{pmatrix},
\qquad
B=\begin{pmatrix}3&4\\0&5\end{pmatrix}.
\]

\item \textbf{Problem 12}\par
If
\[
A^{-1}=\begin{pmatrix}2&-1\\0&3\end{pmatrix}
\quad\text{and}\quad
B^{-1}=\begin{pmatrix}1&0\\4&2\end{pmatrix},
\]
find \((AB)^{-1}\).

\item \textbf{Problem 13}\par
Prove by direct multiplication that
\[
(A^{\mathsf T})^{-1}=(A^{-1})^{\mathsf T}
\]
for
\[
A=\begin{pmatrix}1&2\\3&5\end{pmatrix}.
\]

\item \textbf{Problem 14}\par
Find the inverse of the elementary matrix
\[
E=\begin{pmatrix}
1&0&0\\
0&1&0\\
-4&0&1
\end{pmatrix}
\]
and explain which row operation \(E\) represents.

\item \textbf{Problem 15}\par
A proposed inverse of
\[
A=\begin{pmatrix}2&-1\\1&0\end{pmatrix}
\]
is
\[
C=\begin{pmatrix}0&1\\-1&2\end{pmatrix}.
\]
Check whether the proposal is correct.

\item \textbf{Problem 16}\par
Explain why a matrix with two identical rows cannot have an inverse. Use both a determinant argument and an information-loss interpretation.

\item \textbf{Problem 17}\par
The transformation
\[
y=
\begin{pmatrix}2&1\\1&1\end{pmatrix}x
\]
maps an unknown vector \(x\) to
\[
y=\begin{pmatrix}8\\5\end{pmatrix}.
\]
Recover \(x\).

\item \textbf{Problem 18}\par
Choose the most efficient method for finding the inverse of each matrix and justify the choice:
\[
A=\begin{pmatrix}5&0\\0&-2\end{pmatrix},
\quad
B=\begin{pmatrix}1&2\\3&7\end{pmatrix},
\quad
C=\begin{pmatrix}1&1&0\\0&1&1\\0&0&2\end{pmatrix}.
\]
Then compute the inverses.

\item \textbf{Problem 19}\par
Let
\[
A=\begin{pmatrix}
1&a\\
2&2a
\end{pmatrix}.
\]
Explain, without attempting Gauss--Jordan elimination, why \(A^{-1}\) does not exist for any real \(a\).

\item \textbf{Problem 20}\par
For
\[
A=\begin{pmatrix}
1&2\\
3&4
\end{pmatrix},
\]
find \(A^{-1}\), solve \(Ax=b\) for
\[
b=\begin{pmatrix}5\\11\end{pmatrix},
\]
and verify the solution in the original equation. Finally, explain what the inverse accomplishes conceptually.
\end{enumerate}
